\documentclass[10pt,a4paper]{article}

\usepackage[T1]{fontenc}
\usepackage[utf8]{inputenc}
\usepackage{lmodern}
\usepackage[a4paper, top=1.8cm, bottom=2cm, left=1.5cm, right=1.5cm]{geometry}
\usepackage{amsmath,amssymb,mathtools}
\usepackage{xcolor}
\usepackage{listings}
\usepackage[most]{tcolorbox}
\usepackage{microtype}
\usepackage{needspace}
\usepackage{parskip}
\usepackage{fancyhdr}
\usepackage{hyperref}
\usepackage{multicol}
% sansmath not needed

\setlength{\parskip}{3pt}
\setlength{\parindent}{0pt}

% ─── Colors ───────────────────────────────────────────────────────
\definecolor{codegreen}{rgb}{0.0,0.48,0.18}
\definecolor{codegray}{rgb}{0.50,0.50,0.50}
\definecolor{codepurple}{rgb}{0.38,0.0,0.58}
\definecolor{codeblue}{rgb}{0.0,0.28,0.68}
\definecolor{codeorange}{rgb}{0.70,0.35,0.0}
\definecolor{backcolour}{rgb}{0.965,0.965,0.975}
\definecolor{mathbg}{rgb}{0.92,0.95,1.0}
\definecolor{notebg}{rgb}{1.0,0.97,0.87}
\definecolor{sectionbg}{rgb}{0.16,0.26,0.52}
\definecolor{subsecbg}{rgb}{0.38,0.50,0.75}
\definecolor{bordergray}{rgb}{0.68,0.68,0.74}
\definecolor{noteborder}{rgb}{0.80,0.65,0.25}

% ─── Listings ─────────────────────────────────────────────────────
\lstdefinestyle{gurobipy}{
  language=Python,
  backgroundcolor=\color{backcolour},
  commentstyle=\color{codegray}\itshape,
  keywordstyle=\color{codepurple}\bfseries,
  stringstyle=\color{codegreen},
  basicstyle=\ttfamily\small,
  breakatwhitespace=false,
  breaklines=true,
  keepspaces=true,
  showspaces=false,
  showstringspaces=false,
  showtabs=false,
  tabsize=2,
  frame=none,
  xleftmargin=4pt,
  xrightmargin=4pt,
  aboveskip=0pt,
  belowskip=0pt,
  morekeywords={
    GRB, quicksum, addVar, addVars, addConstr, addConstrs,
    addRange, addQConstr,
    addGenConstrIndicator, addGenConstrAbs,
    addGenConstrMax, addGenConstrMin,
    addSOS, setObjectiveN, setObjective, setParam,
    optimize, update, Model, tupledict, multidict,
    LinExpr, QuadExpr, GurobiError,
    INFINITY, BINARY, CONTINUOUS, INTEGER,
    MINIMIZE, MAXIMIZE, OPTIMAL,
  },
}

% ─── tcolorbox ────────────────────────────────────────────────────
\tcbset{
  mathbox/.style={
    colback=mathbg,
    colframe=codeblue!35,
    boxrule=0.5pt,
    arc=3pt,
    left=6pt, right=6pt, top=5pt, bottom=5pt,
  },
  notebox/.style={
    colback=notebg,
    colframe=noteborder,
    boxrule=0.4pt,
    arc=3pt,
    left=5pt, right=5pt, top=4pt, bottom=4pt,
    fontupper=\footnotesize\itshape,
  },
}

% ─── Lengths ──────────────────────────────────────────────────────
\newlength{\mathcolw}
\setlength{\mathcolw}{0.295\textwidth}
\newlength{\codecolw}
\setlength{\codecolw}{0.665\textwidth}

% ─── Helpers ──────────────────────────────────────────────────────
\newcommand{\secheader}[2]{%
  \vspace{10pt}%
  \needspace{6cm}%
  \phantomsection\hypertarget{sec#1}{}%
  \noindent%
  \begin{tcolorbox}[
    colback=sectionbg, colupper=white,
    colframe=sectionbg, sharp corners,
    left=8pt, right=8pt, top=5pt, bottom=5pt,
    fontupper=\large\bfseries\sffamily,
  ]
  Secci\'on #1:\enspace #2
  \end{tcolorbox}%
  \vspace{2pt}%
}

\newcommand{\entryheader}[2][]{%
  \needspace{5cm}%
  \vspace{6pt}%
  \ifx\relax#1\relax\else\phantomsection\hypertarget{#1}{}\fi%
  \noindent%
  \begin{tcolorbox}[
    colback=subsecbg!18,
    colframe=subsecbg!45,
    sharp corners, boxrule=0.4pt,
    left=5pt, right=5pt, top=2pt, bottom=2pt,
    fontupper=\small\bfseries\sffamily,
  ]
  #2
  \end{tcolorbox}%
  \vspace{2pt}%
}

% shorthand for inline code
\newcommand{\py}[1]{\texttt{\small #1}}

% ─── Header / Footer ──────────────────────────────────────────────
\pagestyle{fancy}
\fancyhf{}
\fancyhead[L]{{\small\sffamily\textcolor{sectionbg}{\textbf{Gurobipy} --- Programaci\'on Matem\'atica en Python}}}
\fancyhead[R]{{\small\sffamily\textcolor{codegray}{Gu\'ia v1.0}}}
\fancyfoot[C]{\small\thepage}
\renewcommand{\headrulewidth}{0.4pt}
\setlength{\headheight}{14pt}

\hypersetup{colorlinks=true, linkcolor=sectionbg, urlcolor=codeblue}

% ══════════════════════════════════════════════════════════════════
\begin{document}
% ══════════════════════════════════════════════════════════════════

% ─── Title ────────────────────────────────────────────────────────
\begin{tcolorbox}[
        colback=sectionbg, colupper=white,
        colframe=sectionbg, sharp corners,
        left=10pt, right=10pt, top=8pt, bottom=8pt,
    ]
    {\LARGE\bfseries\sffamily Hands-on Gurobipy}\\[4pt]
    {\normalsize\sffamily Autor: unknown
        \hfill\textcolor{white!70!sectionbg}{\today}}
\end{tcolorbox}

\vspace{4pt}
\noindent{\footnotesize\textit{%
        Convenciones:\enspace
        $I, J, K$ conjuntos de \'indices;\enspace
        $i, j, k$ \'indices;\enspace
        $a, b, c$ par\'ametros;\enspace
        $M$ constante Big-M;\enspace
        \py{m} $=$ instancia de \py{gb.Model()};\enspace
        \py{x, y, z} variables gurobipy.%
    }}
\vspace{4pt}

% ─── Índice ───────────────────────────────────────────────────────
\begin{tcolorbox}[
    colback=sectionbg!8, colframe=sectionbg!35,
    sharp corners, boxrule=0.4pt,
    left=8pt, right=8pt, top=5pt, bottom=5pt,
]
{\small\bfseries\sffamily\textcolor{sectionbg}{\'{I}ndice de contenidos}}\\[1pt]
\begin{multicols}{2}
\setlength{\columnseprule}{0.3pt}
\footnotesize\setlength{\parskip}{1pt}
\hyperlink{sec1}{\textbf{\S1\enspace Setup del modelo}}\\
\quad\hyperlink{e11}{1.1 Crear modelo e importar librer\'ias}\\
\quad\hyperlink{e12}{1.2 Variable continua con cotas}\\
\quad\hyperlink{e13}{1.3 Variable entera}\\
\quad\hyperlink{e14}{1.4 Variable binaria}\\
\quad\hyperlink{e15}{1.5 Variables indexadas --- Loops}\\
\quad\hyperlink{e16}{1.6 Variables indexadas --- dict comp.}\\
\quad\hyperlink{e17}{1.7 Variables indexadas --- addVars}\\[3pt]
\hyperlink{sec2}{\textbf{\S2\enspace Restricciones b\'asicas}}\\
\quad\hyperlink{e21}{2.1 Igualdad}\\
\quad\hyperlink{e22}{2.2 Desigualdad $\leq$ ponderada}\\
\quad\hyperlink{e23}{2.3 Desigualdad $\geq$ con filtro}\\
\quad\hyperlink{e24}{2.4 Restricciones indexadas --- Loops}\\
\quad\hyperlink{e25}{2.5 Restricciones indexadas --- addConstrs}\\
\quad\hyperlink{e26}{2.6 Rango (restricci\'on doble)}\\[3pt]
\hyperlink{sec3}{\textbf{\S3\enspace Sumas y quicksum}}\\
\quad\hyperlink{e31}{3.1 quicksum vs sum de Python}\\
\quad\hyperlink{e32}{3.2 Objetivo lineal}\\
\quad\hyperlink{e33}{3.3 Suma doble (multi-\'indice)}\\
\quad\hyperlink{e34}{3.4 Objetivo con suma doble}\\[3pt]
\hyperlink{sec4}{\textbf{\S4\enspace Variables binarias y Big-M}}\\
\quad\hyperlink{e41}{4.1 Activaci\'on con Big-M}\\
\quad\hyperlink{e42}{4.2 Equivalencia $y_i=1 \Leftrightarrow x_i>0$}\\
\quad\hyperlink{e43}{4.3 addGenConstrIndicator}\\[3pt]
\hyperlink{sec5}{\textbf{\S5\enspace Restricciones multi-\'indice}}\\
\quad\hyperlink{e51}{5.1 Para todo par $(i,j)$}\\
\quad\hyperlink{e52}{5.2 Suma por fila y columna}\\
\quad\hyperlink{e53}{5.3 Doble suma con filtro}\\[3pt]
\hyperlink{sec6}{\textbf{\S6\enspace Restricciones de red}}\\
\quad\hyperlink{e61}{6.1 Conservaci\'on de flujo}\\
\quad\hyperlink{e62}{6.2 Capacidad de arco binaria}\\
\quad\hyperlink{e63}{6.3 Asignaci\'on: uno activo}\\[3pt]
\hyperlink{sec7}{\textbf{\S7\enspace Linealizaci\'on}}\\
\quad\hyperlink{e71}{7.1 Valor absoluto $t=|x|$}\\
\quad\hyperlink{e72}{7.2 M\'aximo $t=\max_i x_i$}\\
\quad\hyperlink{e73}{7.3 M\'inimo $t=\min_i x_i$}\\[3pt]
\hyperlink{sec8}{\textbf{\S8\enspace Avanzado}}\\
\quad\hyperlink{e81}{8.1 Par\'ametros y lectura de resultados}\\[3pt]
\hyperlink{sec9}{\textbf{\S9\enspace Post-optimizaci\'on}}\\
\quad\hyperlink{e91}{9.1 Precios sombra (dual values)}\\
\quad\hyperlink{e92}{9.2 Holguras (slack/surplus)}\\
\quad\hyperlink{e93}{9.3 Costos reducidos}\\
\quad\hyperlink{e94}{9.4 Ranging del RHS}\\
\quad\hyperlink{e95}{9.5 Ranging coef.\ objetivo}
\end{multicols}
\end{tcolorbox}

\vspace{4pt}

% ══════════════════════════════════════════════════════════════════
\secheader{1}{Setup del modelo}
% ══════════════════════════════════════════════════════════════════

\entryheader[e11]{1.1\quad Crear modelo e importar librer\'ias}
\noindent
\begin{minipage}[t]{\mathcolw}
    \begin{tcolorbox}[mathbox]
        $\min\; f(x)$\\[4pt]
        modelo vac\'io,\\ sin variables\\ ni restricciones
    \end{tcolorbox}
\end{minipage}\hspace{0.02\textwidth}%
\begin{minipage}[t]{\codecolw}
    \begin{lstlisting}[style=gurobipy]
import gurobipy as gb
from gurobipy import GRB, quicksum

m = gb.Model("mi_modelo")     # instancia de modelo
# Opcional para controlar la salida y el tiempo de optimizacion
m.setParam("OutputFlag", 0)   # silenciar salida en consola
m.setParam("TimeLimit", 300)  # limite de tiempo en segundos
    \end{lstlisting}
\end{minipage}

\vspace{4pt}

\entryheader[e12]{1.2\quad Variable continua con cotas}
\noindent
\begin{minipage}[t]{\mathcolw}
    \begin{tcolorbox}[mathbox]
        $x \in [lb,\; ub] \subset \mathbb{R}$
    \end{tcolorbox}
    \vspace{2pt}
    \begin{tcolorbox}[notebox]
        Sin \py{ub}, la variable es ilimitada superiormente (\py{GRB.INFINITY}). \py{name} aparece en logs y en archivos \py{.lp} para debug.
    \end{tcolorbox}
\end{minipage}\hspace{0.02\textwidth}%
\begin{minipage}[t]{\codecolw}
    \begin{lstlisting}[style=gurobipy]
x = m.addVar(
    vtype=GRB.CONTINUOUS,
    lb=0,    # cota inferior (default: 0)
    ub=100,  # cota superior (default: GRB.INFINITY)
    name="x"
)
    \end{lstlisting}
\end{minipage}

\vspace{4pt}

\entryheader[e13]{1.3\quad Variable entera}
\noindent
\begin{minipage}[t]{\mathcolw}
    \begin{tcolorbox}[mathbox]
        $z \in \mathbb{Z}_{\geq 0}$
    \end{tcolorbox}
    \vspace{2pt}
    \begin{tcolorbox}[notebox]
        Se puede combinar con \py{ub} para limitar el rango entero, ejemplo \py{ub=50}.
    \end{tcolorbox}
\end{minipage}\hspace{0.02\textwidth}%
\begin{minipage}[t]{\codecolw}
    \begin{lstlisting}[style=gurobipy]
z = m.addVar(
    vtype=GRB.INTEGER,
    lb=0,               # cota inferior (default: 0)
    ub=GRB.INFINITY,    # opcional, es el default para enteras
    name="z"
)
\end{lstlisting}
\end{minipage}

\vspace{4pt}

\entryheader[e14]{1.4\quad Variable binaria (caso particular de entera con $lb=0$ y $ub=1$)}
\noindent
\begin{minipage}[t]{\mathcolw}
    \begin{tcolorbox}[mathbox]
        $y \in \{0,\; 1\}$
    \end{tcolorbox}
    \vspace{2pt}
    \begin{tcolorbox}[notebox]
        Es equivalente a declarar la variable como entera y limitar su rango a $[0,1]$.\\
        Si se declara como binaria, no es necesario declarar cotas.
    \end{tcolorbox}
\end{minipage}\hspace{0.02\textwidth}%
\begin{minipage}[t]{\codecolw}
    \begin{lstlisting}[style=gurobipy]
y = m.addVar(vtype=GRB.BINARY, name="y")
\end{lstlisting}
\end{minipage}

\vspace{4pt}

\entryheader[e15]{1.5\quad Variables indexadas --- Loops anidados}
\noindent
\begin{minipage}[t]{\mathcolw}
    \begin{tcolorbox}[mathbox]
        $x_{ij} \in [0,1]$\\[4pt]
        $\forall\; i \in I,\; j \in J$
    \end{tcolorbox}
    \vspace{2pt}
    \begin{tcolorbox}[notebox]
        El doble loop es la forma m\'as directa de crear variables indexadas. Sin embargo, para modelos con muchos \'indices, el c\'odigo se vuelve verboso y propenso a errores.
    \end{tcolorbox}
\end{minipage}\hspace{0.02\textwidth}%
\begin{minipage}[t]{\codecolw}
    \begin{lstlisting}[style=gurobipy]
I = ["A", "B", "C"]
J = [1, 2, 3]
X = {}                  # diccionario vacio para guardar las variables
for i in I:             # para cada i en I
    for j in J:         # para cada j en J
        X[i, j] = m.addVar(
            vtype=GRB.CONTINUOUS,
            lb=0, ub=1,
            name=f"X_{i}_{j}"
        )
\end{lstlisting}
\end{minipage}

\vspace{4pt}

\entryheader[e16]{1.6\quad Variables indexadas --- dict comprehension}
\noindent
\begin{minipage}[t]{\mathcolw}
    \begin{tcolorbox}[mathbox]
        $x_{ij} \in [0,1]$\\[4pt]
        $\forall\; i \in I,\; j \in J$
    \end{tcolorbox}
    \vspace{2pt}
    \begin{tcolorbox}[notebox]
        Clave tupla: \py{X["A",1]} $\equiv$ \py{X[("A",1)]}. El dict comprehension es una forma compacta de crear variables indexadas sin escribir bucles anidados. Ideal para modelos con muchos índices.
    \end{tcolorbox}
\end{minipage}\hspace{0.02\textwidth}%
\begin{minipage}[t]{\codecolw}
    \begin{lstlisting}[style=gurobipy]
I = ["A", "B", "C"]
J = [1, 2, 3]

X = {
    (i, j): m.addVar(vtype=GRB.CONTINUOUS,
                     lb=0, ub=1,
                     name=f"X_{i}_{j}")
    for i in I
    for j in J
}
\end{lstlisting}
\end{minipage}

\vspace{4pt}

\entryheader[e17]{1.7\quad Variables indexadas --- \py{addVars} y \py{multidict}}
\noindent
\begin{minipage}[t]{\mathcolw}
    \begin{tcolorbox}[mathbox]
        $x_{ij},\; c_{ij},\; u_{ij}$\\[4pt]
        desde datos estructurados
    \end{tcolorbox}
    \vspace{1pt}
    \begin{tcolorbox}[notebox]
        \py{multidict} devuelve claves + dicts en una sola l\'inea, es decir, \py{arcs} $=$ lista de claves.
    \end{tcolorbox}
    \vspace{1pt}
    \begin{tcolorbox}[notebox]
        \py{addVars} devuelve \py{tupledict} con m\'etodo \py{.sum("A","*")} para sumas parciales sin escribir \py{quicksum}.\\
        Por cada valor en \py{arcs}, se crea una variable con el mismo nombre, ejemplo \py{x["A","B"]}.
    \end{tcolorbox}
\end{minipage}\hspace{0.02\textwidth}%
\begin{minipage}[t]{\codecolw}
    \begin{lstlisting}[style=gurobipy]
arcs, cost, cap = gb.multidict({
    ("A","B"): [2.5, 10],
    ("A","C"): [1.8,  5],
    ("B","C"): [3.0,  8],
})

# addVars devuelve tupledict
x = m.addVars(arcs, vtype=GRB.CONTINUOUS,
              lb=0, name="x")

# suma parcial sobre todo j con origen "A"
expr = x.sum("A", "*")
\end{lstlisting}
\end{minipage}

\vspace{4pt}

% ══════════════════════════════════════════════════════════════════
\secheader{2}{Restricciones b\'asicas}
% ══════════════════════════════════════════════════════════════════

\entryheader[e21]{2.1\quad Igualdad}
\noindent
\begin{minipage}[t]{\mathcolw}
    \begin{tcolorbox}[mathbox]
        $\displaystyle\sum_{i \in I} x_i = b$
    \end{tcolorbox}
    \begin{tcolorbox}[notebox]
        La igualdad se modela con \py{==} en \py{addConstr}.\\
        En modelos grandes, es recomendable usar \py{quicksum} para construir la expresi\'on lineal de manera eficiente.
    \end{tcolorbox}
\end{minipage}\hspace{0.02\textwidth}%
\begin{minipage}[t]{\codecolw}
    \begin{lstlisting}[style=gurobipy]
b = 1.0
m.addConstr(
    quicksum(x[i] for i in I) == b,
    name="igualdad"
)
\end{lstlisting}
\end{minipage}

\vspace{4pt}

\entryheader[e22]{2.2\quad Desigualdad $\leq$ con suma ponderada por $a_i$}
\noindent
\begin{minipage}[t]{\mathcolw}
    \begin{tcolorbox}[mathbox]
        $\displaystyle\sum_{i \in I} a_i\, x_i \leq b$
    \end{tcolorbox}
\end{minipage}\hspace{0.02\textwidth}%
\begin{minipage}[t]{\codecolw}
    \begin{lstlisting}[style=gurobipy]
a = {"A": 3, "B": 2, "C": 5}
b = 10

m.addConstr(
    quicksum(a[i] * x[i] for i in I) <= b,
    name="capacidad"
)
\end{lstlisting}
\end{minipage}

\vspace{4pt}

\entryheader[e23]{2.3\quad Desigualdad $\geq$ --- con restriccion de elementos en $I$}
\noindent
\begin{minipage}[t]{\mathcolw}
    \begin{tcolorbox}[mathbox]
        $\displaystyle\sum_{i \in I \setminus \{B\}} x_i \geq 1$
    \end{tcolorbox}
\end{minipage}\hspace{0.02\textwidth}%
\begin{minipage}[t]{\codecolw}
    \begin{lstlisting}[style=gurobipy]
m.addConstr(
    quicksum(x[i] for i in I if i != "B") >= 1,
    name="cobertura"
)
\end{lstlisting}
\end{minipage}

\vspace{4pt}

\entryheader[e24]{2.4\quad Restricciones indexadas con elementos excluidos de $I$ --- Loops}
\noindent
\begin{minipage}[t]{\mathcolw}
    \begin{tcolorbox}[mathbox]
        $\displaystyle\sum_{\substack{j \in J \\ i \neq j}} x_{ij} \leq u_i, \quad \forall i \in I$
    \end{tcolorbox}
\end{minipage}\hspace{0.02\textwidth}%
\begin{minipage}[t]{\codecolw}
    \begin{lstlisting}[style=gurobipy]
u = {"A": 2, "B": 3, "C": 1}
for i in I:
    m.addConstr(
        quicksum(X[i, j] for j in J if j != i) <= u[i],
        name=f"capacity_{i}"
    )
\end{lstlisting}
\end{minipage}

\vspace{4pt}

\entryheader[e25]{2.5\quad Restricciones indexadas --- Vectorizadas con \py{addConstrs} y \py{quicksum}}
\noindent
\begin{minipage}[t]{\mathcolw}
    \begin{tcolorbox}[mathbox]
        $\displaystyle\sum_{\substack{j \in J \\ i \neq j}} x_{ij} \leq u_i, \quad \forall i \in I$
    \end{tcolorbox}
    \vspace{2pt}
    \begin{tcolorbox}[notebox]
        \py{addConstrs} permite crear múltiples restricciones de manera vectorizada, evitando bucles explícitos.
    \end{tcolorbox}
\end{minipage}\hspace{0.02\textwidth}%
\begin{minipage}[t]{\codecolw}
    \begin{lstlisting}[style=gurobipy]
u = {"A": 2, "B": 3, "C": 1}\
m.addConstrs(
    (quicksum(X[i, j] for j in J if j != i) <= u[i]
     for i in I),
    name="capacity"
)  # crea: capacity[A], capacity[B], capacity[C]
\end{lstlisting}
\end{minipage}

\vspace{4pt}

\entryheader[e26]{2.6\quad Rango (restricci\'on doble)}
\noindent
\begin{minipage}[t]{\mathcolw}
    \begin{tcolorbox}[mathbox]
        $L \leq \displaystyle\sum_{i \in I} x_i \leq U$
    \end{tcolorbox}
    \vspace{2pt}
    \begin{tcolorbox}[notebox]
        \py{addRange} crea 1 fila interna. Dos \py{addConstr} crean 2 filas pero son m\'as legibles en el log.
    \end{tcolorbox}
\end{minipage}\hspace{0.02\textwidth}%
\begin{minipage}[t]{\codecolw}
    \begin{lstlisting}[style=gurobipy]
expr = quicksum(x[i] for i in I)

# opcion 1: una sola restriccion interna
m.addRange(expr, 2, 5, name="rango")

# opcion 2: equivalente, dos restricciones
m.addConstr(expr >= 2, name="rango_lb")
m.addConstr(expr <= 5, name="rango_ub")
\end{lstlisting}
\end{minipage}

\vspace{4pt}

% ══════════════════════════════════════════════════════════════════
\secheader{3}{Sumas y \texttt{quicksum}}
% ══════════════════════════════════════════════════════════════════

\entryheader[e31]{3.1\quad \py{quicksum} vs \py{sum} de Python}
\noindent
\begin{minipage}[t]{\mathcolw}
    \begin{tcolorbox}[mathbox]
        $\displaystyle\sum_{i \in I} c_i\, x_i$
    \end{tcolorbox}
    \vspace{2pt}
    \begin{tcolorbox}[notebox]
        \py{quicksum} construye la \py{LinExpr} sin objetos intermedios. Con $|I|=10^4$, \py{sum()} puede ser hasta $100\times$ m\'as lento porque crea expresiones temporales en cada \py{+}.
    \end{tcolorbox}
\end{minipage}\hspace{0.02\textwidth}%
\begin{minipage}[t]{\codecolw}
    \begin{lstlisting}[style=gurobipy]
c = {"A": 2.0, "B": 3.5, "C": 1.0}

# CORRECTO: eficiente con gurobipy
expr = quicksum(c[i] * x[i] for i in I)

# EVITAR en modelos grandes (crea LinExpr temporales)
expr_lenta = sum(c[i] * x[i] for i in I)
\end{lstlisting}
\end{minipage}

\vspace{4pt}

\entryheader[e32]{3.2\quad Objetivo lineal}
\noindent
\begin{minipage}[t]{\mathcolw}
    \begin{tcolorbox}[mathbox]
        $\displaystyle\min_{x}\;\sum_{i \in I} c_i\, x_i$
    \end{tcolorbox}
\end{minipage}\hspace{0.02\textwidth}%
\begin{minipage}[t]{\codecolw}
    \begin{lstlisting}[style=gurobipy]
m.setObjective(
    quicksum(c[i] * x[i] for i in I),
    GRB.MINIMIZE   # o GRB.MAXIMIZE
)
\end{lstlisting}
\end{minipage}

\vspace{4pt}

\entryheader[e33]{3.3\quad Suma doble (multi-\'indice)}
\noindent
\begin{minipage}[t]{\mathcolw}
    \begin{tcolorbox}[mathbox]
        $\displaystyle\sum_{i \in I}\sum_{j \in J} c_{ij}\, x_{ij}$
    \end{tcolorbox}
    \vspace{2pt}
    \begin{tcolorbox}[notebox]
        \py{cost[i,j]} $\equiv$ \py{cost[(i,j)]}.
    \end{tcolorbox}
\end{minipage}\hspace{0.02\textwidth}%
\begin{minipage}[t]{\codecolw}
    \begin{lstlisting}[style=gurobipy]
cost = {("A",1): 2.0, ("A",2): 1.5,
        ("B",1): 3.0, ("B",2): 2.5}

obj = quicksum(
    cost[i, j] * X[i, j]
    for i in I
    for j in J
)
m.setObjective(obj, GRB.MINIMIZE)
\end{lstlisting}
\end{minipage}

\vspace{4pt}

\entryheader[e34]{3.4\quad Objetivos con suma doble (multi-\'indice)}
\noindent
\begin{minipage}[t]{\mathcolw}
    \begin{tcolorbox}[mathbox]
        $\displaystyle\sum_{i \in I}\sum_{j \in J} c_{ij}\, x_{ij} + \sum_{i \in I} d_i\, y_i$
    \end{tcolorbox}
    \vspace{2pt}
    \begin{tcolorbox}[notebox]
        Se pueden combinar sumas de diferentes \'indices en un mismo objetivo. Solo es necesario construir la expresi\'on completa con \py{quicksum} antes de llamar a \py{setObjective}.
    \end{tcolorbox}
\end{minipage}\hspace{0.02\textwidth}%
\begin{minipage}[t]{\codecolw}
    \begin{lstlisting}[style=gurobipy]
cost = {("A",1): 2.0, ("A",2): 1.5,
        ("B",1): 3.0, ("B",2): 2.5}
d = {"A": 1.0, "B": 0.5, "C": 2.0}
obj = quicksum(
    cost[i, j] * X[i, j]
    for i in I
    for j in J
) + quicksum(
    d[i] * y[i]
    for i in I
)
m.setObjective(obj, GRB.MINIMIZE)
\end{lstlisting}
\end{minipage}

\vspace{4pt}
% ══════════════════════════════════════════════════════════════════
\secheader{4}{Variables binarias y Big-M}
% ══════════════════════════════════════════════════════════════════

\entryheader[e41]{4.1\quad Activaci\'on: $y_i = 1 \Rightarrow x_i \leq M$ y $y_i = 0 \Rightarrow x_i = 0$}
\noindent
\begin{minipage}[t]{\mathcolw}
    \begin{tcolorbox}[mathbox]
        $x_i \leq M \cdot y_i$\\[6pt]
        $\forall\; i \in I$
    \end{tcolorbox}
    \vspace{2pt}
    \begin{tcolorbox}[notebox]
        $M$ debe ser la cota superior \emph{ajustada} de $x_i$.\\
        Ojo: un $M$ grande debilita la relajaci\'on LP y enlentece el Branch \& Bound.
    \end{tcolorbox}
\end{minipage}\hspace{0.02\textwidth}%
\begin{minipage}[t]{\codecolw}
    \begin{lstlisting}[style=gurobipy]
M = 100  # cota superior real de x[i]

for i in I:
    m.addConstr(
        x[i] <= M * y[i],
        name=f"bigM_{i}"
    )
# y[i]=0 -> x[i]=0 ;  y[i]=1 -> x[i] libre hasta M
\end{lstlisting}
\end{minipage}

\vspace{4pt}

\entryheader[e42]{4.2\quad Equivalencia: $y_i = 1 \Leftrightarrow x_i > 0$}
\noindent
\begin{minipage}[t]{\mathcolw}
    \begin{tcolorbox}[mathbox]
        $x_i \leq M \cdot y_i$\\[2pt]
        $x_i \geq \varepsilon \cdot y_i$\\[6pt]
        $\forall\; i \in I$
    \end{tcolorbox}
    \vspace{2pt}
    \begin{tcolorbox}[notebox]
        $\varepsilon$ es un umbral peque\~no ($10^{-3}$) que fuerza $y_i = 1$ cuando $x_i > 0$. Solo v\'alido si $x_i$ es continua.
    \end{tcolorbox}
\end{minipage}\hspace{0.02\textwidth}%
\begin{minipage}[t]{\codecolw}
    \begin{lstlisting}[style=gurobipy]
eps = 1e-3

for i in I:
    m.addConstr(x[i] <= M * y[i],
                name=f"bigM_ub_{i}")
    m.addConstr(x[i] >= eps * y[i],
                name=f"bigM_lb_{i}")
\end{lstlisting}
\end{minipage}

\vspace{4pt}

\entryheader[e43]{4.3\quad \py{addGenConstrIndicator} --- implicaci\'on limpia}
\noindent
\begin{minipage}[t]{\mathcolw}
    \begin{tcolorbox}[mathbox]
        $y_i = 1 \Rightarrow$\\[2pt]
        $\quad\displaystyle\sum_{j \in J} x_{ij} \leq C_i$\\[6pt]
        $\forall\; i \in I$
    \end{tcolorbox}
    \vspace{2pt}
    \begin{tcolorbox}[notebox]
        M\'as legible que Big-M. Gurobi lo convierte internamente a Big-M durante el preproceso. No garantiza ventaja sobre un Big-M bien formulado.
    \end{tcolorbox}
\end{minipage}\hspace{0.02\textwidth}%
\begin{minipage}[t]{\codecolw}
    \begin{lstlisting}[style=gurobipy]
C = {"A": 5, "B": 3}

for i in I:
    m.addGenConstrIndicator(
        y[i],   # binaria "switch"
        True,   # valor del switch: True=1, False=0
        quicksum(X[i, j] for j in J) <= C[i],
        name=f"indicator_{i}"
    )
\end{lstlisting}
\end{minipage}

\vspace{4pt}

% ══════════════════════════════════════════════════════════════════
\secheader{5}{Restricciones multi-\'indice}
% ══════════════════════════════════════════════════════════════════

\entryheader[e51]{5.1\quad Para todo par $(i,\,j)$}
\noindent
\begin{minipage}[t]{\mathcolw}
    \begin{tcolorbox}[mathbox]
        $x_{ij} \leq y_i$\\[6pt]
        $\forall\; i \in I,\; j \in J$
    \end{tcolorbox}
\end{minipage}\hspace{0.02\textwidth}%
\begin{minipage}[t]{\codecolw}
    \begin{lstlisting}[style=gurobipy]
for i in I:
    for j in J:
        m.addConstr(
            X[i, j] <= y[i],
            name=f"link_{i}_{j}"
        )
\end{lstlisting}
\end{minipage}

\vspace{4pt}

\entryheader[e52]{5.2\quad Suma por fila y por columna}
\noindent
\begin{minipage}[t]{\mathcolw}
    \begin{tcolorbox}[mathbox]
        $\displaystyle\sum_{j \in J} x_{ij} = 1 \quad \forall\, i$\\[6pt]
        $\displaystyle\sum_{i \in I} x_{ij} \leq 1 \quad \forall\, j$
    \end{tcolorbox}
\end{minipage}\hspace{0.02\textwidth}%
\begin{minipage}[t]{\codecolw}
    \begin{lstlisting}[style=gurobipy]
# restriccion por fila (para cada i)
for i in I:
    m.addConstr(
        quicksum(X[i, j] for j in J) == 1,
        name=f"fila_{i}"
    )

# restriccion por columna (para cada j)
for j in J:
    m.addConstr(
        quicksum(X[i, j] for i in I) <= 1,
        name=f"col_{j}"
    )
\end{lstlisting}
\end{minipage}

\vspace{4pt}

\entryheader[e53]{5.3\quad Doble suma con filtro}
\noindent
\begin{minipage}[t]{\mathcolw}
    \begin{tcolorbox}[mathbox]
        $\displaystyle\sum_{k \in K} f_{ikt}\, x_{ikt}$\\[4pt]
        $\leq \displaystyle\sum_{\substack{q \in Q\\ c_{iq}>0}} c_{iq}\, z_{iqt}$\\[6pt]
        $\forall\; i \in I,\; t \in T$
    \end{tcolorbox}
    \vspace{2pt}
    \begin{tcolorbox}[notebox]
        El filtro \py{if cap[i,q] > 0} evita t\'erminos nulos que engordan el modelo sin aportar restricci\'on real.
    \end{tcolorbox}
\end{minipage}\hspace{0.02\textwidth}%
\begin{minipage}[t]{\codecolw}
    \begin{lstlisting}[style=gurobipy]
for t in T:
    for i in I:
        m.addConstr(
            quicksum(
                f[i, k, t] * X[i, k, t]
                for k in K)
            <=
            quicksum(
                cap[i, q] * Z[i, q, t]
                for q in Q
                if cap[i, q] > 0),  # filtrar ceros
            name=f"cap_{i}_{t}"
        )
\end{lstlisting}
\end{minipage}

\vspace{4pt}

% ══════════════════════════════════════════════════════════════════
\secheader{6}{Restricciones de red}
% ══════════════════════════════════════════════════════════════════

\entryheader[e61]{6.1\quad Conservaci\'on de flujo}
\noindent
\begin{minipage}[t]{\mathcolw}
    \begin{tcolorbox}[mathbox]
        $\displaystyle\sum_{\substack{j:\,(i,j)\in A}} x_{ij}$\\[4pt]
        $-\displaystyle\sum_{\substack{j:\,(j,i)\in A}} x_{ji} = b_i$\\[6pt]
        $\forall\; i \in V$
    \end{tcolorbox}
    \vspace{2pt}
    \begin{tcolorbox}[notebox]
        $b_i < 0$: fuente;\enspace $b_i > 0$: sumidero;\enspace $b_i = 0$: tr\'ansito. Los arcos se filtran din\'amicamente desde la lista \py{arcs}.
    \end{tcolorbox}
\end{minipage}\hspace{0.02\textwidth}%
\begin{minipage}[t]{\codecolw}
    \begin{lstlisting}[style=gurobipy]
nodes = ["s", "1", "2", "t"]
arcs  = [("s","1"), ("s","2"),
         ("1","t"), ("2","t")]
b     = {"s": -1, "1": 0, "2": 0, "t": 1}

x = {a: m.addVar(lb=0, name=f"x_{a}")
     for a in arcs}

for i in nodes:
    salida  = quicksum(x[u,v] for (u,v) in arcs
                       if u == i)
    entrada = quicksum(x[u,v] for (u,v) in arcs
                       if v == i)
    m.addConstr(salida - entrada == b[i],
                name=f"flujo_{i}")
\end{lstlisting}
\end{minipage}

\vspace{4pt}

\entryheader[e62]{6.2\quad Capacidad de arco con apertura binaria}
\noindent
\begin{minipage}[t]{\mathcolw}
    \begin{tcolorbox}[mathbox]
        $0 \leq x_{ij} \leq u_{ij} \cdot y_{ij}$\\[6pt]
        $\forall\; (i,j) \in A$
    \end{tcolorbox}
\end{minipage}\hspace{0.02\textwidth}%
\begin{minipage}[t]{\codecolw}
    \begin{lstlisting}[style=gurobipy]
u = {("s","1"): 10, ("s","2"): 5,
     ("1","t"):  8, ("2","t"): 6}

y = {a: m.addVar(vtype=GRB.BINARY,
                 name=f"y_{a}")
     for a in arcs}

for a in arcs:
    m.addConstr(x[a] <= u[a] * y[a],
                name=f"cap_{a}")
\end{lstlisting}
\end{minipage}

\vspace{4pt}

\entryheader[e63]{6.3\quad Asignaci\'on: exactamente uno activo}
\noindent
\begin{minipage}[t]{\mathcolw}
    \begin{tcolorbox}[mathbox]
        $\displaystyle\sum_{i \in I} y_i = 1$
    \end{tcolorbox}
\end{minipage}\hspace{0.02\textwidth}%
\begin{minipage}[t]{\codecolw}
    \begin{lstlisting}[style=gurobipy]
m.addConstr(
    quicksum(y[i] for i in I) == 1,
    name="exactamente_uno"
)
\end{lstlisting}
\end{minipage}

\vspace{4pt}

% ══════════════════════════════════════════════════════════════════
\secheader{7}{Linealizaci\'on de no-linealidades}
% ══════════════════════════════════════════════════════════════════

\entryheader[e71]{7.1\quad Valor absoluto: $t = |x|$}
\noindent
\begin{minipage}[t]{\mathcolw}
    \begin{tcolorbox}[mathbox]
        $t \geq x$\\
        $t \geq -x$\\
        $t \geq 0$\\[4pt]
        $\Leftrightarrow\; t = |x|$\\(si se minimiza $t$)
    \end{tcolorbox}
    \vspace{2pt}
    \begin{tcolorbox}[notebox]
        Las dos formas son equivalentes. \py{addGenConstrAbs} es m\'as conciso; Gurobi lo convierte a las dos desigualdades durante el preproceso.
    \end{tcolorbox}
\end{minipage}\hspace{0.02\textwidth}%
\begin{minipage}[t]{\codecolw}
    \begin{lstlisting}[style=gurobipy]
x = m.addVar(lb=-GRB.INFINITY, name="x")
t = m.addVar(lb=0, name="t_abs")

# Forma manual (dos restricciones lineales)
m.addConstr(t >= x,  name="abs_pos")
m.addConstr(t >= -x, name="abs_neg")

# Forma directa con restriccion generalizada
m.addGenConstrAbs(t, x, name="abs_x")
# minimizar t en objetivo -> t = |x|
\end{lstlisting}
\end{minipage}

\vspace{4pt}

\entryheader[e72]{7.2\quad M\'aximo: $t = \max_i x_i$}
\noindent
\begin{minipage}[t]{\mathcolw}
    \begin{tcolorbox}[mathbox]
        $t = \displaystyle\max_{i \in I} x_i$
    \end{tcolorbox}
    \vspace{2pt}
    \begin{tcolorbox}[notebox]
        Solo cotas $t \geq x_i$ es suficiente si $t$ entra en el objetivo con coeficiente $> 0$ (minimizaci\'on). Para igualdad exacta en otros contextos usar \py{addGenConstrMax}.
    \end{tcolorbox}
\end{minipage}\hspace{0.02\textwidth}%
\begin{minipage}[t]{\codecolw}
    \begin{lstlisting}[style=gurobipy]
t_max = m.addVar(lb=-GRB.INFINITY, name="t_max")

# Opcion 1: solo cotas inferiores (si minimizando t)
for i in I:
    m.addConstr(t_max >= x[i],
                name=f"max_lb_{i}")

# Opcion 2: igualdad exacta garantizada
m.addGenConstrMax(
    t_max,
    [x[i] for i in I],
    constant=-GRB.INFINITY,  # sin constante adicional
    name="max_x"
)
\end{lstlisting}
\end{minipage}

\vspace{4pt}

\entryheader[e73]{7.3\quad M\'inimo: $t = \min_i x_i$}
\noindent
\begin{minipage}[t]{\mathcolw}
    \begin{tcolorbox}[mathbox]
        $t = \displaystyle\min_{i \in I} x_i$
    \end{tcolorbox}
\end{minipage}\hspace{0.02\textwidth}%
\begin{minipage}[t]{\codecolw}
    \begin{lstlisting}[style=gurobipy]
t_min = m.addVar(lb=-GRB.INFINITY, name="t_min")

m.addGenConstrMin(
    t_min,
    [x[i] for i in I],
    constant=GRB.INFINITY,  # sin constante adicional
    name="min_x"
)
\end{lstlisting}
\end{minipage}

\vspace{4pt}

% ══════════════════════════════════════════════════════════════════
\secheader{8}{Avanzado}
% ══════════════════════════════════════════════════════════════════

\entryheader[e81]{8.1\quad Par\'ametros del solver y lectura de resultados}
\noindent
\begin{minipage}[t]{\mathcolw}
    \begin{tcolorbox}[mathbox]
        Configuraci\'on\\ del solver\\[4pt]
        Lectura de\\resultados tras\\ \py{m.optimize()}
    \end{tcolorbox}
\end{minipage}\hspace{0.02\textwidth}%
\begin{minipage}[t]{\codecolw}
    \begin{lstlisting}[style=gurobipy]
m.setParam("MIPGap",    0.01) # gap optimalidad 1%
m.setParam("TimeLimit", 3600) # 1 hora maximo
m.setParam("Threads",   4)    # hilos en paralelo
m.setParam("MIPFocus",  1)    # buscar factible rapido
m.setParam("Presolve",  2)    # presolve agresivo
m.setParam("OutputFlag", 1)   # activar log en consola

m.optimize()

# Leer resultados
if m.Status == GRB.OPTIMAL:
    print(f"ObjVal:   {m.ObjVal:.4f}")
    print(f"ObjBound: {m.ObjBound:.4f}")
    print(f"MIPGap:   {100*m.MIPGap:.2f}%")
    print(f"Runtime:  {m.Runtime:.1f}s")
    sol = {i: x[i].X for i in I}  # .X = valor optimo
\end{lstlisting}
\end{minipage}


% ══════════════════════════════════════════════════════════════════
\secheader{9}{An\'alisis de sensibilidad post-optimizaci\'on}
% ══════════════════════════════════════════════════════════════════

\entryheader[e91]{9.1\quad Precios sombra (\textit{dual values})}
\noindent
\begin{minipage}[t]{\mathcolw}
    \begin{tcolorbox}[mathbox]
        $\lambda_i = \dfrac{\partial z^*}{\partial b_i}$\\[6pt]
        $\Delta z^* \approx \lambda_i \cdot \Delta b_i$\\[4pt]
        $\forall\; i \in \text{restricciones}$
    \end{tcolorbox}
    \vspace{2pt}
    \begin{tcolorbox}[notebox]
        Solo LP (no MIP directo). En min: $\lambda_i > 0$ $\Rightarrow$ relajar $b_i$ reduce $z^*$. Restricci\'on ``cara'': activa con $\lambda_i \neq 0$.
    \end{tcolorbox}
\end{minipage}\hspace{0.02\textwidth}%
\begin{minipage}[t]{\codecolw}
    \begin{lstlisting}[style=gurobipy]
m.optimize()  # LP requerido para analisis de sensibilidad

for c in m.getConstrs():
    print(f"{c.ConstrName:15s}  Pi = {c.Pi:+.4f}")

# Pi = -2.5 (min): relajar b_i en 1 reduce z* en 2.5
c_cap = m.getConstrByName("capacidad")
print(f"precio sombra: {c_cap.Pi:.4f}")
    \end{lstlisting}
\end{minipage}

\vspace{4pt}

\entryheader[e92]{9.2\quad Holguras (\textit{slack\,/\,surplus})}
\noindent
\begin{minipage}[t]{\mathcolw}
    \begin{tcolorbox}[mathbox]
        $s_i = b_i - \displaystyle\sum_j a_{ij}\,x_j^*$\\[4pt]
        $s_i = 0$:\enspace restricci\'on activa\\[2pt]
        $s_i > 0$:\enspace restricci\'on inactiva
    \end{tcolorbox}
    \vspace{2pt}
    \begin{tcolorbox}[notebox]
        Restricci\'on inactiva $\Rightarrow\;\lambda_i = 0$. Para $\geq$:\; $s_i = \text{lhs} - b_i$. Disponible en LP y MIP.
    \end{tcolorbox}
\end{minipage}\hspace{0.02\textwidth}%
\begin{minipage}[t]{\codecolw}
    \begin{lstlisting}[style=gurobipy]
for c in m.getConstrs():
    activa = abs(c.Slack) < 1e-6
    estado = "activa" if activa else "inactiva"
    print(f"{c.ConstrName}: slack={c.Slack:.4f}  [{estado}]")
    \end{lstlisting}
\end{minipage}

\vspace{4pt}

\entryheader[e93]{9.3\quad Costos reducidos (\textit{reduced costs})}
\noindent
\begin{minipage}[t]{\mathcolw}
    \begin{tcolorbox}[mathbox]
        $\bar{c}_j = c_j - \mathbf{c}_B^\top \mathbf{B}^{-1} \mathbf{a}_j$\\[6pt]
        Var.\ b\'asica:\enspace $\bar{c}_j = 0$\\[2pt]
        No b\'asica (min):\enspace $\bar{c}_j \geq 0$
    \end{tcolorbox}
    \vspace{2pt}
    \begin{tcolorbox}[notebox]
        $|\bar{c}_j|$ = cu\'anto debe cambiar $c_j$ para que $x_j$ entre a la base. Si $x_j^* > 0$, entonces RC $= 0$.
    \end{tcolorbox}
\end{minipage}\hspace{0.02\textwidth}%
\begin{minipage}[t]{\codecolw}
    \begin{lstlisting}[style=gurobipy]
for v in m.getVars():
    estado = "basica" if abs(v.RC) < 1e-6 else "no basica"
    print(
        f"{v.VarName:10s}  x*={v.X:.4f}"
        f"  RC={v.RC:+.4f}  [{estado}]"
    )
    \end{lstlisting}
\end{minipage}

\vspace{4pt}

\entryheader[e94]{9.4\quad Ranging del RHS}
\noindent
\begin{minipage}[t]{\mathcolw}
    \begin{tcolorbox}[mathbox]
        $b_i \in [\text{Low}_i,\;\text{Up}_i]$\\[6pt]
        $\lambda_i$ constante en ese\\ intervalo (misma base)\\[4pt]
        Fuera del rango:\\ la base cambia
    \end{tcolorbox}
    \vspace{2pt}
    \begin{tcolorbox}[notebox]
        \py{SARHSLow/Up}: l\'imites de $b_i$ donde $\lambda_i$ sigue siendo v\'alido. Si $b_i$ sale del rango, la base cambia.
    \end{tcolorbox}
\end{minipage}\hspace{0.02\textwidth}%
\begin{minipage}[t]{\codecolw}
    \begin{lstlisting}[style=gurobipy]
for c in m.getConstrs():
    print(
        f"{c.ConstrName:15s}"
        f"  b_i in [{c.SARHSLow:.3f}, {c.SARHSUp:.3f}]"
        f"  Pi={c.Pi:.4f}"
    )
# [Low, Up]: intervalo donde lambda_i permanece constante
    \end{lstlisting}
\end{minipage}

\vspace{4pt}

\entryheader[e95]{9.5\quad Ranging de coeficientes objetivo}
\noindent
\begin{minipage}[t]{\mathcolw}
    \begin{tcolorbox}[mathbox]
        $c_j \in [\text{Low}_j,\;\text{Up}_j]$\\[6pt]
        Base \'optima no cambia\\ en ese intervalo
    \end{tcolorbox}
    \vspace{2pt}
    \begin{tcolorbox}[notebox]
        Fuera del rango, otra variable puede entrar a la base y $x^*$ cambia. \'Util para evaluar robustez frente a cambios en costos.
    \end{tcolorbox}
\end{minipage}\hspace{0.02\textwidth}%
\begin{minipage}[t]{\codecolw}
    \begin{lstlisting}[style=gurobipy]
for v in m.getVars():
    print(
        f"{v.VarName:10s}"
        f"  c_j={v.Obj:.3f}"
        f"  in [{v.SAObjLow:.3f}, {v.SAObjUp:.3f}]"
    )
# SAObjLow/Up: rango de c_j donde la base actual
# sigue siendo optima (misma solucion x*)
    \end{lstlisting}
\end{minipage}

\vspace{14pt}

% ─── Footer note ──────────────────────────────────────────────────
\begin{tcolorbox}[
        colback=sectionbg!8, colframe=sectionbg!30,
        sharp corners, boxrule=0.4pt,
        left=8pt, right=8pt, top=4pt, bottom=4pt,
    ]
    \footnotesize
    \textbf{Referencia:} Gurobi Optimizer Reference Manual v10.0\enspace
\end{tcolorbox}

\end{document}
